\section*{Technical Experience}

\textbf{Bayesian Inference Pipelines}
\begin{itemize}
    \item Designed and developed \href{https://github.com/seap-udea/PRisma}{PRisma}, a modular Bayesian inference framework for exoplanetary research, integrating forward modeling, KDE-based likelihood estimation, and nested sampling (\texttt{dynesty}). Implemented parallel execution on HPC clusters, reproducible workflows, and version-controlled inference campaigns, enabling large-scale analyses and supporting ongoing first-author research.
\end{itemize}
\textit{Python, Bayesian inference, Nested sampling (dynesty), Parallel computing, HPC, Scientific computing, Reproducible research}

\vspace{6pt}
\textbf{Applied Data Mining \& Machine Learning}
\begin{itemize}
    \item Applied end-to-end data science workflows to astronomical datasets, including SQL/ADQL querying, data preprocessing, visualization, feature engineering, and supervised and unsupervised machine learning. Developed predictive and clustering models using the Python scientific ecosystem as part of research-oriented projects.
\end{itemize}
\textit{SQL, ADQL, scikit-learn, Data analysis, Classification, Clustering, Machine learning, Data visualization}

\vspace{6pt}
\textbf{Scientific Software Engineering}
\begin{itemize}
    \item Co-developed and maintained \href{https://github.com/seap-udea/pryngles}{Pryngles}, an open-source Python package for photometric simulations of planetary systems. Designed and evolved its public API, implemented testing and documentation infrastructure, and applied modern software engineering practices including modular architecture, collaborative Git workflows, reproducible development, and long-term maintainability.
\end{itemize}
\textit{Python, Software architecture, API design, Open-source development, Git / GitHub, Unit testing, Documentation, Reproducible research}
